%============================================================================
% prelude.tex
%   - titlepage
%   - table of contents,and list of figures
%   - abstract
%============================================================================
\clearpage%
\pagenumbering{roman}  % This makes the page numbers Roman (i, ii, etc)
%--------------------------------------------------------------------%
% TITLE PAGE
%   - define \title{} \author{} \date{}

\title{{ OPTICAL CHARACTER RECOGNITION}}

\author{Tawade Shubham Ratnoji (15ET1105) \\ Singh Vishal Kumar  (15ET1138) \\ Sonawane Riddesh Rajesaheb(15ET1043)}

\date{November 2018}

\department{Department of Electronics and Telecommunication Engineering}

%    - Set the guide's name
\setguide{Mr. Jaswantsing Rajput}
%    - Set the coguide's name (if you have one)
%\setcoguide{Prof. Name of the Co-supervisor}

%   - once the above are defined, use \maketitle to generate the titlepage
\maketitle

%--------------------------------------------------------------------
\cleardoublepage \nonumber
\begin{center}
	\vspace*{-1in}
\includegraphics{raitlogo}\\
Ramrao Adik Education Society's\\
\large {\textbf{Ramrao Adik Institute of Technology}}\\
\small(Affiliated to the University of Mumbai)\\
Dr. D. Y. Patil Vidyanagar, Sector 7, Nerul, Navi Mumbai 400
706.\\\vspace{0.4in}
\textbf{\begin{Huge}Certificate of Approval\end{Huge}}
\end{center}
%\vspace{0.in}
\begin{center}
\large This is to certify that, the Project -I report entitled
\vspace{0.15in}\\\textbf {\lq\lq OPTICAL CHARACTER RECOGNITION\rq\rq} \\
\vspace{0.1in}is a bonafide work done by\vspace{0.15in}
\\\textbf {Tawade Shubham Ratnoji (15ET1105) \\ Singh Vishal Kumar  (15ET1138) \\ Sonawane Riddesh Rajesaheb(15ET1043) } \\\vspace{0.1in} and is submitted in the partial
fulfillment of
the requirement for the degree of \vspace{0.15in} \\\textbf{Bachelor of Engineering}\\\textbf{(Electronics and Telecommunication Engineering)} 
\\to the \\\textbf{University of Mumbai}. \\
\includegraphics[width=0.3\textwidth]{Figures/university.png}
 \vspace{0.2in}
\begin{tabular}{ccc}

      \rule{4.0cm}{1sp} & \hspace{0.7in}                & \hspace{0.3in} \rule{4.0cm}{1sp} \\ \vspace{0.65in}
       Examiner & \hspace{1in}             & \hspace{0.3in} Supervisor \\
      \rule{4.0cm}{1sp} & \hspace{0.5in} \rule{4.0cm}{1sp}                  &\hspace{0.45in} \rule{4.0cm}{1sp} \\
      Project Coordinator & \hspace{0.5in} Head of Department                &\hspace{0.65in} Principal \\\\
        %&\rule{5cm}{1sp}& \\
%       &Principal&
    \end{tabular}
\end{center}

%----------------------------------------------------------------------------------
\cleardoublepage \nonumber
\begin{center}
	\vspace*{-1in}
	\includegraphics{raitlogo}\\
	Ramrao Adik Education Society's\\
	\large {\textbf{Ramrao Adik Institute of Technology}}\\
	\small(Affiliated to the University of Mumbai)\\
	Dr. D. Y. Patil Vidyanagar,Sector 7, Nerul, Navi Mumbai 400
	706.\\\vspace{0.4in}
	\textbf{\begin{Huge}Declaration\end{Huge}}
\end{center}
%\vspace{0.in}

	We wish to state that work embodied in this dissertation entitled ``\textbf{OPTICAL CHARACTER RECOGNITION}'' has been carried out under the guidance
	of Mr. Jaswantsing Rajput at Department of Electronics and Telecommunication Engineering, Ramrao Adik Institute of Technology during 2018-2019. 
	\par
	We declare that the work being presented forms our own contribution and has not been submitted for any other Degree or Diploma of any University/Institute. Wherever references have been made to previous works of others, it has been clearly indicated. We understand that any violation of the above will be cause for disciplinary action by the Institute and can also evoke penal action from the sources which have thus not been properly cited or from whom proper permission has not been taken when needed.
	\\
	\vspace{3cm}
	\\
	\begin{center}
		\begin{tabular}{cc}
			\rule{5.0cm}{1sp} & \hspace{4cm} \rule{5.0cm}{1sp} \\ 
			Tawade Shubham Ratnoji & \hspace{4cm} Singh Vishal Kumar \\
			& \\ &\\ &\\ &\\
			\rule{5.0cm}{1sp} &  \\
			Sonawane Riddesh Rajasaheb\\
		\end{tabular}
	\end{center}

%-------------------------------------------------------------------------------
\cleardoublepage
\begin{acknowledgments}
We are very glad to thank to our project guide \textbf{Mr. Jaswantsing Rajput} and our \textbf{H.O.D Dr. M.D.Patil} for their encouragement and tremendous guidance. Here it is very special thank feelings for our colleagues for
their support and help. We have been fortunate to have received many useful suggestions
from our colleagues which have greatly improved the clarity of our report. At the end
special thanks to our \textbf{Principal Dr. Ramesh Vasappanavara}. We would like to appreciate
suggestions and criticisms about the report from the readers.
\end{acknowledgments}

% ABSTRACT
\begin{abstract}
 \paragraph{}In India, more than 310 million people use Devanagari script for documentation. There has been a significant improvement in the research related to the recognition of printed as well as handwritten Devanagari text in the past few years. All feature-extraction techniques as well as training, classification and matching techniques useful for the recognition are used.
\paragraph{}In any OCR or classification system extracting discriminating feature is most important and crucial step for its success. Accuracy of such system often depends on the good feature representation. In this project, we have used histogram of individual characters and GLCM(Gray Level Co-occurence Matrix) for feature extraction. In this part, first image is preprocessed to remove noise, converted to binary image, resized to fixed size and then convert to gray scale image using mask operation, it blurs the edges of the images.  Finally network weight parameters are fine tuned by supervised back propagation learning to improve the overall recognition performance. 
\paragraph{} Now-a-days there are many new methodologies required for the increasing needs in newly emerging areas, with these methodologies there are many techniques are present for the character recognition of handprint Devnagari, Bangla, Tamil, China etc. But it has not been used practically. So in this project, we have used a Minimum distance classifier technique for OCR System of printed as well as scanned newsprint Marathi script.            

\end{abstract}

%--------------------------------------------------------------------%
% CONTENTS, TABLES, FIGURES
 \tableofcontents

%--------------------------------------------------------------------%
%  Single counter for theorems and theorem-like environments:
\newtheorem{theorem}{Theorem}[chapter]
\newtheorem{assertion}[theorem]{Assertion}
\newtheorem{claim}[theorem]{Claim}
\newtheorem{conjecture}[theorem]{Conjecture}
\newtheorem{corollary}[theorem]{Corollary}
\newtheorem{definition}[theorem]{Definition}
\newtheorem{example}[theorem]{Example}
\newtheorem{figger}[theorem]{Figure}
\newtheorem{lemma}[theorem]{Lemma}
\newtheorem{prop}[theorem]{Proposition}
\newtheorem{remark}[theorem]{Remark}

%--------------------------------------------------------------------%
% Make the page numbers Arabic (1, 2, etc)
\cleardoublepage%
\pagenumbering{arabic} 